% Extracted from p02-moe-routing
\subsection{Mixture-of-Experts Routing at Sub-100M Scale}
\label{ch-moe}

\paragraph{Background}

\paragraph{Top-$k$ routing.} For token representation $x$, a router
computes logits $g = W_r x \in \mathbb{R}^E$ and probabilities
$p = \mathrm{softmax}(g)$. With top-$k$ routing, the token is
dispatched to the $k$ experts with largest $p_i$, weighted by their
(renormalised) probabilities. The Switch Transformer
\citep{fedus2022switch} popularised $k=1$ with a load-balancing
auxiliary loss; we keep $k$ general and specialise where needed.

\paragraph{The auxiliary loss.} Let $f_i$ be the fraction of tokens in
a batch dispatched to expert $i$ and $P_i$ the mean routing probability
assigned to expert $i$. The Switch load-balancing loss is
\begin{equation}
  \mathcal{L}_{\mathrm{aux}} = \alpha \, E \sum_{i=1}^{E} f_i P_i ,
  \label{eq:aux}
\end{equation}
with coefficient $\alpha > 0$. We derive its behaviour in
Section~\ref{sec:aux}.

\paragraph{Routing as estimation: the small-batch problem}
\label{sec:tokens}

The router is a classifier whose ``labels'' are implicit: a token
should go to the expert that will process it best. Its gradient at
each step is estimated from the tokens routed to each expert. With
$N$ tokens per optimisation step and $E$ experts under top-1 routing,
token counts follow a multinomial distribution, so for expert $i$ with
routing share $q_i$,
\begin{equation}
  \mathbb{E}[n_i] = N q_i, \qquad
  \mathrm{Var}(n_i) = N q_i (1 - q_i).
\end{equation}
Under uniform routing ($q_i = 1/E$), the relative standard deviation
of expert $i$'s token count is
\begin{equation}
  \frac{\sqrt{\mathrm{Var}(n_i)}}{\mathbb{E}[n_i]}
  = \sqrt{\frac{E - 1}{N}} \approx \frac{1}{\sqrt{N/E}}.
\end{equation}
The router's gradient for expert $i$'s parameters is an average over
$n_i$ tokens; its noise therefore scales roughly as $1/\sqrt{N/E}$
(the approximation is coarse at $E=8$, off by about 6\%). The
quantity $N/E$---tokens per expert per step---is the router's
effective sample size.

For our 44M configuration, $N \approx 120{,}004{,}608 / 4{,}850
\approx 2.47 \times 10^4$ tokens per step and $E = 8$, giving
$N/E \approx 3{,}090$ tokens per expert per step under uniform
routing, with relative fluctuation around $1.8\%$. That is
comfortable \emph{at the step level}. But two caveats bite at small
scale. First, if gradient accumulation is used, the per-microbatch
$N$ can be an order of magnitude smaller, pushing relative
fluctuations past $10\%$. Second, and more fundamentally, small
models train on far fewer tokens in total---120M here against the
trillions of large runs---so the router gets far fewer updates over
training to average out the noise or to recover from an early
bad lock-in. The router is thus solving a bandit problem with few
pulls per arm per round and few rounds overall: the classic
conditions for premature exploitation.

\paragraph{The rich-get-richer map}
\label{sec:feedback}

Beyond noise, there is a deterministic feedback. Let
$f^{(t)} \in \Delta^{E-1}$ be the routed fractions at step $t$ and
$\ell^{(t)}$ the corresponding mean router logits, with
$f^{(t)} = \mathrm{softmax}(\ell^{(t)})$. Expert $i$'s parameters move
in proportion to the tokens it sees, so to first order its logit
grows in proportion to its routed fraction:
\begin{equation}
  \ell^{(t+1)} = \ell^{(t)} + \alpha f^{(t)}, \qquad \alpha > 0,
\end{equation}
where $\alpha$ absorbs the learning rate, the tokens per step, and
the mean per-token gradient alignment. The uniform point
$\ell = c\mathbf{1}$ ($f_i = 1/E$) is a fixed point. Linearising
around it in the zero-sum gauge ($\sum_i \epsilon_i = 0$),
\begin{equation}
  f_i \approx \frac{1}{E} + \frac{\epsilon_i}{E}, \qquad
  \epsilon' = \frac{\alpha}{E}\,\epsilon .
\end{equation}
Perturbations therefore grow geometrically whenever
\begin{equation}
  \alpha > E .
  \label{eq:instability}
\end{equation}
Since $\alpha = \eta N \bar{g}$ (learning rate $\times$ tokens per
step $\times$ mean per-token logit gradient), condition
\eqref{eq:instability} is easier to satisfy when per-token gradients
are large and coherent---precisely the early-training regime where
some experts randomly fit the data better than others. The map is
deliberately simple, but it isolates the mechanism: routing
concentration is \emph{self-amplifying}, and the uniform distribution
is an unstable equilibrium once feedback exceeds the threshold.
Noise (Section~\ref{sec:tokens}) decides \emph{which} experts win;
the feedback decides that \emph{someone} wins everything.

\paragraph{The load-balancing loss, derived}
\label{sec:aux}

The auxiliary loss \eqref{eq:aux} penalises the correlation between
dispatched fractions $f_i$ and mean probabilities $P_i$. In the
idealised case $f_i = P_i$,
\begin{equation}
  \sum_{i=1}^E f_i P_i = \sum_{i=1}^E f_i^2 \geq \frac{1}{E},
\end{equation}
by convexity, with equality iff routing is uniform. The factor $E$
normalises the minimum to 1, so $\mathcal{L}_{\mathrm{aux}} \geq
\alpha$ with equality at perfect balance. The loss is minimised by
spreading tokens evenly---but note it does not distinguish
\emph{which} tokens go where, only the marginal distribution. It is
a regulariser on the router's marginals, not a teaching signal for
routing quality.

\paragraph{Choosing $\alpha$.} The auxiliary term contributes
$O(\alpha)$ per token to the total loss. If $\alpha$ is too large,
the router optimises balance at the expense of the language-modelling
objective (uniform routing with random experts); if too small, the
feedback of Section~\ref{sec:feedback} dominates. A practical,
scale-free heuristic: choose $\alpha$ so that
$\mathcal{L}_{\mathrm{aux}}$ is roughly $1$--$10\%$ of the LM loss at
initialisation, then verify the diagnostics of
Section~\ref{sec:diagnostics} rather than tuning blindly. The Switch
Transformer's $\alpha = 0.01$ \citep{fedus2022switch} sits in this
range for typical setups, but the right value moves with $E$ and
batch size, so it should be re-derived per configuration, not copied.

\paragraph{Capacity, dropping, and utilisation}
\label{sec:capacity}

In practice each expert accepts at most
$C = \gamma \cdot kN/E$ tokens per batch ($\gamma \geq 1$ the
capacity factor); overflow tokens are dropped (passed through
unprocessed). Define utilisation $U_i = \min(1, n_i / C)$.
Dropping interacts badly with collapse: once an expert is
under-utilised, its effective batch shrinks further, its gradients
get noisier, and recovery becomes less likely. A capacity factor
$\gamma$ moderately above 1 absorbs the multinomial fluctuation
derived in Section~\ref{sec:tokens} without masking systematic
imbalance---dropping should be rare under healthy routing and common
only under collapse, which makes the drop rate itself a diagnostic.

\paragraph{Proposed diagnostics}
\label{sec:diagnostics}

We propose logging, per MoE layer and per evaluation window:

\begin{enumerate}
  \item \textbf{Routing entropy.}
    $H = -\sum_i \bar{q}_i \log \bar{q}_i$ over corpus-level routing
    shares $\bar{q}_i$, normalised by $\log E$. Healthy: near 1.
    Collapsing: trending toward 0.
  \item \textbf{Gini coefficient} of per-expert token counts.
    A single number in $[0,1]$; 0 is perfect balance.
  \item \textbf{Dead-expert fraction:} share of experts with
    $\bar{q}_i$ below $1/(10E)$ over the window. Nonzero is an alarm.
  \item \textbf{Drop rate:} fraction of tokens dropped by capacity
    limits. Spikes indicate the imbalance has become structural.
\end{enumerate}
The point of these metrics is early warning: by the time validation
loss reflects collapse, the router has usually locked in and the
starved experts cannot recover within a small model's token budget.

\paragraph{Discussion}

Applied to our setting---$E=8$ experts, $\approx 2.5 \times 10^4$
tokens per step, 120M training tokens on TinyStories---the analysis
suggests the risk is real but manageable: step-level token counts
per expert are adequate, while microbatch-level noise and the short
total training horizon are the threats. The prescriptions are
unglamorous: keep $\alpha$ honest via the initialisation heuristic,
log the four diagnostics from the first step, and treat any nonzero
dead-expert fraction as a stop-and-fix event rather than a curiosity.
Jitter on router logits (as in Switch) and a router-specific lower
learning rate are compatible additions, but they treat symptoms;
the diagnostics treat the disease by making it visible.

\paragraph{Limitations}

This paper is entirely analytical: we report no new experiments, and
the dynamical map of Section~\ref{sec:feedback} is illustrative, not
a fitted model of any real router. We assume top-$k$ routing with
$k$ small; expert-choice routing and learned temperature schedules
are out of scope. TinyStories is a simple distribution; routing
behaviour on more diverse corpora may differ quantitatively.

\paragraph{Conclusion}

Small MoE models are not small versions of large MoE models where
routing is concerned. The router learns from $N/E$ tokens per step,
and at sub-100M scale that number, together with short training
horizons, makes expert collapse a first-order risk rather than a
curiosity. The instability condition $\alpha > E$, the auxiliary
loss, and the four diagnostics give practitioners a complete,
cheap toolkit: understand the feedback, regularise it, and watch it.
